\documentclass[12pt,a4paper,oneside]{article}

% Pacotes
\usepackage[utf8]{inputenc}
\usepackage[T1]{fontenc}
\usepackage[brazil]{babel}
\usepackage{geometry}
\geometry{a4paper, margin=2.5cm}
\usepackage{indentfirst}
\usepackage{times}
\usepackage{graphicx}
\usepackage{float}
\usepackage{hyperref}
\usepackage{amsmath}
\usepackage{abntex2cite}

% Formatação
\setlength{\parindent}{1.25cm}
\setlength{\parskip}{0.3cm}

\begin{document}

\begin{center}
  \textbf{\large ESCOLA E FACULDADE DE TECNOLOGIA SENAI GASPAR RICARDO JÚNIOR}\\
  \textbf{Desenvolvimento de Sistemas}\\[1cm]
  \textbf{\Large Lucas Miguel Leite}\\[0.5cm]
  \textbf{Professor: Deivison Takatu}\\[2cm]
  \textbf{\LARGE\textbf{Análise de Aplicação de Sistema Digital de Controle Distribuído (SDCD)}}\\[1cm]
  Sorocaba -- 20/03/2026
\end{center}

\vspace{1cm}

\begin{center}
  \textbf{RESUMO}
\end{center}

Este artigo analisa a aplicação de um Sistema Digital de Controle Distribuído (SDCD) em um processo contínuo, com base na arquitetura tecnológica composta por Sensor de Temperatura, ESP32, protocolo MQTT, Servidor, Aplicativo Mobile, Dashboard e interface de Gestão da Fábrica. O estudo detalha o papel de cada componente no sistema, o fluxo de dados desde o nível de campo até o nível de supervisão, as decisões e ações possíveis com base nas informações centralizadas, e as vantagens do uso de sistemas distribuídos em comparação aos modelos centralizados tradicionais. Conclui-se que a arquitetura proposta oferece eliminação de pontos únicos de falha, redução de cabeamento, confiabilidade e escalabilidade, representando uma evolução significativa em relação aos sistemas de controle centralizados.

\vspace{0.5cm}
\textbf{Palavras-chave:} SDCD. Controle Distribuído. MQTT. ESP32. IoT Industrial. SCADA.

\newpage

\section{Introdução}

Os Sistemas Digitais de Controle Distribuído (SDCD) representam uma evolução fundamental no contexto da automação industrial, substituindo os modelos centralizados tradicionais por arquiteturas em que o processamento e a inteligência são distribuídos ao longo da rede, próximos aos processos controlados \cite{moraes2007}. Este artigo analisa a aplicação de um SDCD com base na seguinte arquitetura tecnológica:

\begin{center}
  Sensor de Temperatura $\rightarrow$ ESP32 $\rightarrow$ MQTT $\rightarrow$ Servidor $\rightarrow$ Aplicativo Mobile $\rightarrow$ Dashboard $\rightarrow$ Gestão da Fábrica
\end{center}

\section{Papel de Cada Componente do Sistema}

\subsection{Sensor de Temperatura}

O sensor de temperatura atua no nível de campo. Sua função é monitorar a grandeza física (temperatura) do ambiente ou equipamento industrial e coletar essas informações brutas para o sistema. É o ponto inicial de entrada de dados na cadeia de controle \cite{lara2021}.

\subsection{ESP32 (Controlador Distribuído)}

O ESP32 atua como o controlador inteligente localizado próximo ao processo controlado. Ele recebe os sinais analógicos/digitais do sensor de temperatura, realiza o processamento local das informações e atua como gateway de conectividade. Esta é a essência do controle distribuído: a inteligência reside próximo ao campo, não em um painel central distante \cite{albuquerque2009}.

\subsection{Protocolo MQTT}

O protocolo MQTT (\textit{Message Queuing Telemetry Transport}) representa a camada de comunicação digital do sistema. É o protocolo de rede responsável por empacotar e transmitir os dados lidos pelo ESP32 de maneira leve, rápida e confiável até o servidor central. O MQTT opera no modelo publicador/assinante, ideal para redes industriais com restrição de banda e necessidade de baixa latência.

\subsection{Servidor}

O servidor atua como o núcleo de recepção, armazenamento e roteamento das informações recebidas dos controladores na rede de alta velocidade. É responsável por manter a base de dados histórica e atual que alimenta as interfaces de supervisão.

\subsection{Aplicativo Mobile e Dashboard (Nível de Supervisão)}

O aplicativo mobile e o dashboard operam como a Interface Homem-Máquina (IHM) e sistema SCADA (\textit{Supervisory Control and Data Acquisition}) na sala de controle ou dispositivo remoto. Eles traduzem os dados armazenados no servidor em gráficos, indicadores e alertas visuais compreensíveis para os operadores e gestores.

\section{Fluxo de Dados no Sistema}

O fluxo de dados ocorre de maneira estruturada, partindo do nível de campo até o nível de supervisão, seguindo as etapas \cite{groover2011}:

\begin{enumerate}
  \item \textbf{Coleta:} O sensor captura as variações térmicas do ambiente físico de forma ininterrupta.
  \item \textbf{Processamento Local e Publicação:} O ESP32 lê os dados elétricos, converte em valores digitais compreensíveis (por exemplo, graus Celsius) e utiliza o protocolo MQTT para publicar essas informações na rede.
  \item \textbf{Roteamento e Armazenamento:} O Servidor recebe as mensagens MQTT de temperatura vindas do ESP32 e atualiza sua base de dados.
  \item \textbf{Apresentação e Visualização:} O Aplicativo Mobile e o Dashboard solicitam (ou assinam) esses dados do servidor e atualizam as telas do usuário final. Dessa forma, a gestão da fábrica visualiza o status em tempo real de forma integrada.
\end{enumerate}

\section{Decisões e Ações Baseadas no Dashboard}

A partir do monitoramento centralizado integrado, a gestão pode tomar diversas ações estratégicas e operacionais:

\subsection{Controle de Processos}

É possível identificar desvios na temperatura ideal de produção e enviar um comando remoto de volta (através da rede MQTT) para que o controlador ative um atuador, como o acionamento de um sistema de resfriamento. Isso demonstra a capacidade de controle bidirecional do sistema.

\subsection{Prevenção de Falhas e Manutenção}

A observação de padrões anormais de aquecimento (tendências) pode indicar desgaste ou atrito mecânico de um equipamento, permitindo o agendamento de manutenção preditiva antes que a máquina quebre. Esta abordagem migrada da manutenção preventiva baseada em tempo para a manutenção preditiva baseada em condição representa um ganho significativo em disponibilidade \cite{santos2024}.

\subsection{Segurança e Monitoramento}

É possível configurar alertas visuais e sonoros no aplicativo caso a temperatura atinja níveis críticos, permitindo o desligamento de emergência da linha ou a evacuação da área. Esta funcionalidade é essencial para a segurança do trabalho em ambientes industriais com processos contínuos.

\section{Vantagens do Sistema Distribuído}

A implementação da arquitetura de SDCD proposta traz melhorias significativas em comparação aos sistemas de controle centralizados tradicionais \cite{caiçara2015}:

\subsection{Eliminação de Pontos Únicos de Falha}

O processamento ocorre de forma inteligente e local no ESP32. Se o servidor ou a rede cair, o ESP32 ainda pode estar programado para desligar a máquina automaticamente em caso de superaquecimento extremo, garantindo segurança mesmo sem comunicação com o nível superior.

\subsection{Redução Drástica de Cabeamento}

Como a comunicação via MQTT pode utilizar redes modernas (e muitas vezes sem fio, no caso de microcontroladores como o ESP32), elimina-se o alto custo de materiais e mão de obra para interligar fisicamente todos os sensores do chão de fábrica até um único painel na sala de controle.

\subsection{Confiabilidade e Redundância}

Sistemas distribuídos permitem múltiplos caminhos de comunicação. Se um nó falhar, os demais podem continuar operando, e a topologia de rede pode ser projetada com rotas alternativas para garantir a continuidade do monitoramento \cite{lugli2009}.

\subsection{Escalabilidade}

É muito mais simples adicionar novos sensores e microcontroladores (ESP32) à rede MQTT existente para monitorar outras partes da fábrica, sem a necessidade de paralisar a produção contínua do processo. Cada novo nó é independente e se integra à rede de forma modular.

\section{Conclusão}

A análise da arquitetura de SDCD proposta demonstra que os sistemas de controle distribuído representam uma evolução significativa em relação aos modelos centralizados. A combinação de microcontroladores inteligentes (ESP32), protocolos leves (MQTT) e interfaces de supervisão modernas (dashboards e aplicativos mobile) cria um ecossistema robusto, escalável e seguro para o monitoramento e controle de processos industriais contínuos.

As vantagens --- eliminação de pontos únicos de falha, redução de cabeamento, confiabilidade e escalabilidade --- posicionam essa arquitetura como uma solução viável tanto para indústrias que buscam modernizar sua base instalada quanto para novas instalações que desejam partir diretamente para o paradigma da Indústria 4.0 \cite{schwab2016}.

\bibliographystyle{plain}
\bibliography{referencias}

\end{document}
